% Requires: \usetikzlibrary{positioning,arrows.meta}
\begin{figure}[H]
  \centering
  \begin{tikzpicture}[
      >={Stealth[length=1.8mm]},
      res/.style={draw=black!70,line width=0.4pt,fill=black!6,
        rounded corners=2pt,align=center,inner sep=3.5pt,
        font=\footnotesize},
      note/.style={font=\scriptsize,text=black!55,align=center},
      branch/.style={draw=black,line width=0.5pt,->},
      lead/.style={draw=black!55,line width=0.35pt,dashed}
    ]
    \node[res] (P) at (0,0)
      {$(O,C)$\\
       \textcolor{black!60}{\scriptsize free set $F$,\quad $r=k-|O|$}};
    \node[res] (L) at (-2.5,-2.1)
      {$(O\cup\{e\},C)$\\
       \textcolor{black!60}{\scriptsize open $e$}};
    \node[res] (R) at (2.5,-2.1)
      {$(O,C\cup\{e\})$\\
       \textcolor{black!60}{\scriptsize close $e$}};
    \draw[branch] (P) -- (L);
    \draw[branch] (P) -- (R);
    % P1 acts on the popped node, covering counts against r
    \node[note,anchor=east] (p1) at (-3.1,0.5)
      {P1, covering\\ $\max_l s_I+\max_l s_J>r$};
    \draw[lead] (p1.east) -- (P.west);
    % P2 acts on the popped node, admissible bound
    \node[note,anchor=west] (p2) at (3.1,0.5)
      {P2, bound\\ $LB(O,C)\ge\min(\mathrm{best},B+1)$};
    \draw[lead] (p2.west) -- (P.east);
    % P3 acts only at a leaf
    \node[note] (p3) at (0,-3.0)
      {P3, dominance, only at a leaf, where $F=\varnothing$ and the
       oracle flow leaves some facility of $O$ unused};
  \end{tikzpicture}
  \caption{One branching step of Algorithm~\ref{alg:xpbb}. A node
  carries the forced-open set $O$, the forced-closed set $C$, the free
  set $F$ and the remaining cardinality $r$, and branching on a free
  facility $e$ produces the two children along the solid arrows. The
  dashed leaders attach P1 and P2 to the popped node they test, while
  P3 fires only at a leaf, so the three prunes act exactly on the named
  quantities (Lemma~\ref{lem:cover},
  Proposition~\ref{prop:bbcorrect}).}
  \label{fig:bbtree}
\end{figure}
